% Transcription — fonds Alexandre Grothendieck, shelfmark 140-1, batch 2 (pages 21-36).
% Established from the facsimile archives/batches/140-1/batch-02.pdf (cut from
% https://grothendieck.umontpellier.fr/140-1.pdf, PDF page k = archivists'
% page 20+k, no cover sheet), per the skill transcribe-grothendieck. The folder
% is thirty-six pages; this is the second and last batch (pages 1-20 are batch
% 1, transcribed separately; its notation is kept here).
%
% Pass: Opus 5.5 (claude-opus-5-5), 2026-09-24 — first pass, unchecked against the pages by a human.
% The facsimile and the tiles were prepared beforehand; the archive host was
% not reached from this session.
%
% Pages skipped: 21 (blank sheet), 23, 25, 27 and 36 (typed leaves with
% nothing in his hand). Pages 22 and 24 are written over the verso of a typed
% leaf whose text shows through; only his writing is transcribed. Page 26
% carries a printed address label, page 33 a short note in red ink in another
% hand (dated 27-11-81, recorded once in a \note because it is the only date
% on these leaves), neither transcribed. Pages 29, 33 and 35 are partly
% written upside down; the inverted parts are given in reading order with a
% \note.
%
% His own pagination: none on these leaves. Pages 22 and 24 carry his letters
% « A) » and « B) » (a lemma and the end of its proof); page 31 has « Cas a) »,
% « Cas b) »; page 33 opens with a struck « 4) » and a « 1) ». The section
% headings are the edition's and say so.
%
% What the batch is about. It continues batch 1 on the orbifold (U_{0,3},
% S_3) and the groups Pi^D_{0,3}, SL(2,Z), T_{1,1}, ST_{1,1}. Pp. 22, 24: a
% lemma — S regular integral with Pic(S) = 0 and generic characteristic
% not 2 — giving H^1(U_{0,3}, S_3; G_m) = Z/6Z (Z/3Z in characteristic 2),
% generated by the principal bundle V(O(1))^* on Sigma_J = P(V_J); proof by
% reducing the structure group to mu_6 via the section prod (e_i - e_j) of
% O(6), with the invariants gamma_2, gamma_3 of Sym(V_J)^{S_3}. P. 26: two
% cubes of maps between M_{11}, M_{11}(2), M_{11}(2~), PM_{11}, the stacks
% (U_{0,3}, {1,omega}), (U_{0,3}, D'_3), (U_{0,3}, D_3) (« gal », « gerbe »),
% and the same cube with (script D, Pi_{0,3}), (script D, SL(2,Z)). Pp. 28,
% 30: large (pencil, then ink) diagrams of the extensions Z -> Pi~ -> Pi, the
% groups D_3, D'_3, D''_3, D'''_3 of indices 2 and 3 over Z/12Z, Z/6Z, Z/4Z,
% Z/2Z, and Pi~^D_{0,3}(4) -> Pi^D_{0,3} x Z/4Z. P. 29: sections e_i - e_j of
% O(1), V(omega) = V, and (upside down) the effect of reversing the
% orientation of a combinatorial map on rho_s, rho_f. P. 31 (blue ink): the
% construction of the rho-orbits of a combinatorial map (A, sigma, A', rho),
% finite and infinite cases, and the identification of quadruples with
% triples (A, sigma, rho). Pp. 32, 34: the class of an extension of Z/pZ x
% Z/qZ by Z in Ext^1 = Z/mu Z, then p = 6, q = 4, mu = 12, and the character
% Pi^D_{0,3} -> Z/6Z defined by the Z/6Z-covering PU_{03} = O(1)^*; T_{1,1} =
% Pi^D_{0,3} x_{Z/12Z} Z/4Z. Pp. 33, 35: ST_{11} ~ SL~(2,Z) ~ <rho, sigma |
% rho^3 = sigma^2> ~ Pi^D_{03} x_{Z/6Z} Z ~ SPi^D_{0,3}, and SPi^D_{03}(4)
% with its maps to ST_{11}, T_{11}.
%
% Hardest pages: 31 (fast blue drafting, much connective prose \ill{}), 29
% (mixed orientations, prose at the top largely lost), 28 (a perspective
% drawing of two cubes, crossed by long pencil strokes; laid out on a grid,
% placement partly uncertain and said so).
\documentclass[11pt,a4paper]{article}
% Le préambule commun à toutes les transcriptions du fonds Grothendieck.
%
% Il tient en une idée : **l'incertitude se compose, elle ne se résout pas.**
% Une transcription qui lisse un mot douteux a détruit la seule chose pour
% laquelle on la faisait — un lecteur doit pouvoir savoir, à chaque endroit,
% ce qui était sur le feuillet et ce qui a été ajouté. Les six macros qui
% suivent sont donc l'appareil critique, et non de la décoration.
%
% Les mêmes commandes sont reconnues par `scripts/render.mjs`, qui produit la
% vue de lecture du site. Une macro ajoutée ici doit l'être aussi là-bas,
% faute de quoi le rendu échouera bruyamment — ce qui est voulu : un rendu
% silencieusement faux est pire qu'un rendu absent.

\ProvidesPackage{grothendieck}[2026/08/08 Transcriptions du fonds Grothendieck]

\usepackage{amsmath,amssymb,amsthm}
\usepackage{mathrsfs}
\usepackage{stmaryrd}
% Les diagrammes commutatifs sont partout dans ce fonds : c'est la langue
% même de la géométrie algébrique de Grothendieck, et une transcription qui
% les rendrait en prose ne transcrirait rien.
\usepackage{tikz-cd}
% Français par défaut — c'est la langue des feuillets — mais l'anglais est
% chargé aussi : la traduction et le résumé sont des documents anglais, et la
% typographie française (espaces avant les deux-points) y serait une faute.
% Ces documents basculent par \selectlanguage{english} en tête de corps.
\usepackage[english,french]{babel}
\usepackage{fontspec}
\usepackage{geometry}
\geometry{a4paper,margin=25mm,marginparwidth=28mm}
\usepackage{marginnote}
\usepackage{xcolor}
% ulem, not soul. `soul`'s \st and \ul are built on a character-by-character
% scan that XeLaTeX's font machinery breaks: the document fails with an
% undefined control sequence rather than degrading. ulem does the same two jobs
% and is Unicode-engine safe. `normalem` stops it hijacking \emph, which the
% transcriptions use for Grothendieck's own emphasis.
\usepackage[normalem]{ulem}
\usepackage{hyperref}
\hypersetup{colorlinks=true,linkcolor=black,urlcolor=black,pdfborder={0 0 0}}

% --- Métadonnées du lot -----------------------------------------------------
% Reprises telles quelles par le script de rendu, qui les affiche en tête de
% la vue de lecture. Elles fixent ce que le fichier prétend couvrir : sans
% elles, une transcription flotte sans point d'attache dans un fonds de seize
% mille feuillets.

\newcommand{\folder}[1]{\gdef\grothFolder{#1}}
\newcommand{\batch}[1]{\gdef\grothBatch{#1}}
\newcommand{\leaves}[2]{\gdef\grothFirst{#1}\gdef\grothLast{#2}}
\newcommand{\foldertitle}[1]{\gdef\grothFoldertitle{#1}}
\gdef\grothFolder{?}\gdef\grothBatch{?}\gdef\grothFirst{?}\gdef\grothLast{?}\gdef\grothFoldertitle{}

% --- Le résumé --------------------------------------------------------------
%
% Il ouvre la lecture modernisée, sous le titre « Résumé », et il est composé
% en petit. Ce n'est pas un ornement : c'est le seul passage du document qui
% ne soit pas des mathématiques, et un lecteur doit pouvoir le reconnaître
% avant de l'avoir lu — puis le sauter s'il connaît déjà le sujet, ou s'y
% arrêter s'il ne le connaît pas. Le filet qui le suit marque la reprise du
% corps.

\newenvironment{resume}
  {\small\setlength{\parskip}{0.45em}}
  {\par\medskip\hrule\medskip}

% --- Mots-clés ---------------------------------------------------------------
%
% En anglais, en clôture du résumé de la lecture modernisée : le vocabulaire
% moderne sous lequel chercher ce que les feuillets construisent. Le manifeste
% du site (`npm run manifest`) les extrait pour étiqueter la cote entière —
% cette ligne est la source unique des tags, il n'y a pas de fichier à côté.
\newcommand{\keywords}[1]{\par\smallskip\noindent
  {\footnotesize\textsc{Keywords} --- \textit{#1}}\par}

% --- L'appareil critique ----------------------------------------------------

% Le numéro de feuillet, en marge. C'est l'ancre qui permet de régler un
% désaccord sur une formule en regardant *un* feuillet plutôt qu'une chemise
% de six cents. Le site s'en sert aussi pour tourner les pages du fac-similé
% quand on fait défiler la transcription.
\newcommand{\leaf}[1]{%
  \marginnote{\footnotesize\color{gray!70}\textsf{#1}}%
  \label{leaf:#1}\ignorespaces}

% Illisible : on ne devine pas. Un mot inventé qui se lit comme les autres est
% le pire résultat possible.
\newcommand{\ill}{\textcolor{red!60!black}{[\,\ldots\,]}}

% Lecture douteuse : le mot est donné, et signalé comme incertain.
\newcommand{\uncertain}[1]{\uline{#1}}

% Ajout de l'éditeur — une lettre manquante, un mot sous-entendu.
\newcommand{\add}[1]{\textcolor{blue!50!black}{[#1]}}

% Biffé par Grothendieck lui-même. Ce qu'il a rayé fait partie du document :
% une rature dit souvent où la pensée a changé de direction.
\newcommand{\struck}[1]{\sout{#1}}

% Note du transcripteur, sur l'état matériel du feuillet.
\newcommand{\note}[1]{\par\smallskip{\small\color{gray!60!black}\textit{[#1]}}\par\smallskip}

% Note marginale de Grothendieck, distincte du corps du texte.
\newcommand{\marginal}[1]{\par\smallskip\noindent\hspace{1em}%
  {\small\color{gray!60!black}#1}\par\smallskip}

% --- Titre ------------------------------------------------------------------

\AtBeginDocument{%
  \begin{center}
    {\large\bfseries Fonds Alexandre Grothendieck --- cote n\textsuperscript{o}~\grothFolder}\\[2pt]
    {\itshape\grothFoldertitle}\\[4pt]
    {\small feuillets \grothFirst--\grothLast\ (lot \grothBatch)}\\[2pt]
    {\footnotesize\color{gray!60!black}Transcription établie d'après le fac-similé numérisé
     par l'Université de Montpellier.\\
     Les lectures incertaines sont soulignées, les illisibles notées [\,\ldots\,],
     les ajouts entre crochets.}
  \end{center}
  \vspace{1em}}

\endinput


\folder{140-1}
\batch{2}
\pages{21}{36}
\dating{[à partir de 1978-à partir de 1982]}
\watermark{Édition de démonstration}
\foldertitle{"Longue Marche" [brouillon] : notes manuscrites (s.d.)}

\begin{document}

\section{Le groupe $H^1(U_{0,3}, \mathfrak{S}_3; \mathbb{G}_m)$}

\note{titre de l'édition. Les pages 22 et 24 sont écrites sur le verso d'un
feuillet dactylographié dont le texte transparaît ; seule son écriture est
transcrite. Elles portent, en tête à gauche, « A) » et « B) », entourés d'un
arc}

\page{22}

A) \textbf{Lemme} $S$ schéma \struck{satisfaisant} régulier intègre tel que
$\mathrm{Pic}(S) = 0$, \uncertain{et} car. du point générique $\neq 2$. Alors
\[
H^1(U_{0,3}, \mathfrak{S}_3; \mathbb{G}_m) \simeq \mathbb{Z}/6\mathbb{Z}
\]
(si \struck{résiduelle} \uncertain{car.} \uncertain{générique} $= 2$, on
trouve $\mathbb{Z}/3\mathbb{Z}$)

\note{le « et » de la deuxième ligne est une insertion au-dessus de la
ligne, reliée par une flèche ; le « Alors » est suivi d'un mot biffé}

\marginal{\textbf{Dém.} On utilise la suite exacte
$0 \to H^1(\mathfrak{S}_3, H^0(U_{0,3}, \mathbb{G}_m))$
$\to H^1(U_{0,3}, \mathfrak{S}_3; \mathbb{G}_m)$}

\marginal{$\to \underbrace{H^0(\mathfrak{S}_3, H^1(U_{0,3}; \mathbb{G}_m))}_{0}$ et
$0 \to H^0(S, \mathbb{G}_m) \to H^0(U_{0,3}, \mathbb{G}_m)$
$\to V(J) \to 0$}

\note{note écrite en oblique dans la marge de gauche, sur quatre lignes. Le
terme médian de la première suite est coupé par le bord ; sous
$H^0(\mathfrak{S}_3, \ldots)$ une accolade marquée $0$. Entre
$H^0(U_{0,3}, \ldots)$ et $V(J)$, un signe lu \uncertain{$r_i$}}

…\ dis que si $\Sigma_J$ (le schéma) \struck{\ill{}} est défini comme
$\mathbb{P}(V_J)$, où
\[
0 \to V_J \to \mathcal{O}^J \to \mathcal{O} \to 0 ,
\]
alors le fibré principal
\[
P\Sigma_J = \mathbb{V}(\mathcal{O}(1))^{*}_{\Sigma_J}
\]
est un générateur du $H^1(\Sigma_J^{*}, \mathfrak{S}_J; \mathbb{G}_m)$.

\note{« $\Sigma_J$ (le schéma) » est écrit au-dessus d'un mot biffé.
L'astérisque de $\Sigma^{*}_J$ est récrit sur une rature}

Je vais l'établir \uncertain{en} \uncertain{trouvant} une trivialisation de
\[
\mathcal{O}(1)^{\otimes 6} = \mathcal{O}(6) \simeq V_J^{\otimes 3}(\omega),
\]
\uncertain{définie} \struck{sur} une section \ill{} sur $\Sigma^{*}$. \ill{}
de $\mathcal{O}(6)$ (qui \struck{\ill{} \ill{}}) \struck{\uncertain{prenant}}
\[
\Gamma(\Sigma_J, \mathcal{O}(6)) \simeq \mathrm{Sym}^6 V_J
\]

\note{ligne très surchargée : entre « $\Sigma^{*}$ » et « de
$\mathcal{O}(6)$ », deux lignes de mots insérés et biffés ; seuls quelques
mots se lisent}

d. on prendra
\[
\varphi \in \Gamma(\Sigma_J, \mathcal{O}(6)), \qquad
\varphi = \prod_{\substack{i, j \in J \\ i \neq j}} (e_i - e_j)
= \prod_{i \in J} \underbrace{(e_j - e_k)^2}_{J = \{i, j, k\}}
\]

\note{avant le premier produit, un signe $\pm$ \uncertain{} biffé}

Ainsi, sur $\Sigma^{*}_J$, le groupe structural de $P\Sigma^{*}_J$ se
réduit à : $\mu_6$.

\textbf{NB} \uncertain{Signe} \uncertain{erroné}, …
\[
\struck{\ill{}}\ \mathrm{Sym}^{*}(V_J)^{\mathfrak{S}_3} \simeq
\mathcal{O}[\gamma_2, \gamma_3]/(2\gamma_3^2 - \gamma_2^3)
\]
\[
\begin{cases}
\gamma_2 = \sum_{i \in J} (e_j - e_k)^2 \\
\gamma_3 = \prod_{i \in J} (e_j + e_k - 2e_i)
\end{cases}
\qquad \text{et} \quad \Delta = \gamma_2^{\,3}
\]

\note{les indices de $\gamma_2$, $\gamma_3$ dans le crochet sont récrits ;
le dernier terme du dénominateur et l'égalité $\Delta = \gamma_2^3$ se lisent
mal. Le signe du produit $\gamma_3$ est lu $+$ entre $e_j$ et $e_k$}

\page{24}

B) \struck{Si 2 \uncertain{inversible} \ill{}, on \ill{} \ill{}}

\struck{$\mathrm{Sym}^{*}(V_J)^{\mathfrak{S}_3} \simeq \mathcal{O}[$}

\struck{\ill{} \ill{}}

Comme $\gamma_2$ s'annule sur les \struck{\ill{}} $P_i$, $\gamma_3$ sur les
$Q_i$, il s'ensuit qu'il n'y a pas de section inversible d'un
$\mathcal{O}(n)^{*}$ \struck{\ill{}} sur $\Sigma_J^{*}$ \struck{de degré
$<$} pour $n \uncertain{\leqslant} 6$, donc que l'ordre de
$\mathrm{cl}(\mathcal{O}(1))$ dans
$H^1(\Sigma_J, \mathfrak{S}_J; \mathbb{G}_m)$ est bien $6$, ce qui prouve
que $\mathcal{O}(1)$ est un générateur du $H^1$ \uncertain{cherché}.

\note{« $n \leqslant 6$ » : le signe se lit $\leqslant$, là où l'argument
demande $n < 6$}

\section{Les champs $(U_{0,3}, \mathbb{D}'_3)$ et les revêtements de $M_{11}$}

\note{titre de l'édition. Page 26 : trois diagrammes sans texte. Un trait de
crayon discontinu, fait de courts segments obliques parallèles, traverse les
deux diagrammes du bas sans rien biffer}

\page{26}

\begin{tikzcd}
\mathbb{D}'_3 \arrow[r] & \mathbb{D}_3
\end{tikzcd}

\begin{tikzcd}[column sep=small, row sep=small, nodes={font=\scriptsize}]
PM_{11}(\tilde 2) \simeq PU_{03} \arrow[r, Rightarrow] \arrow[d] & PM_{11}(2) \simeq TU_{03} \arrow[d] \\
U_{0,3} \simeq M_{11}(\tilde 2) \arrow[r] & (U_{0,3}, \{1, \omega\}) \simeq M_{11}(2)
\end{tikzcd}

\note{le carré est marqué « cart ». Au-dessus du premier $\simeq$, un petit
« $2$ » \uncertain{}. Dans le second $M_{11}(\tilde 2)$, l'indice $11$ est
repassé et le $\tilde 2$ porte une barre au-dessus du tilde}

\[
\struck{x \longmapsto x^2 \wedge \omega_0}
\]
\[
PU(TU_{03}, \mathbb{D}_3; \uncertain{\mu_2})
\]

\begin{tikzcd}[column sep=small, row sep=small, nodes={font=\scriptsize}]
M_{11}(\tilde 2) \arrow[rr] \arrow[dr] & & M_{11}(2) \arrow[dl] \arrow[d] \\
 & M_{11} \arrow[d] & M_{11}(2)' \arrow[dl] \\
 & M'_{11} &
\end{tikzcd}

\note{la flèche $M_{11}(2) \to M_{11}(2)'$ est barrée d'un petit trait. Sous
$M_{11}(\tilde 2)$, au départ de la flèche vers $M_{11}$, une étiquette
surchargée où se lit \uncertain{$/\Pi_{0,3}$}}

\begin{tikzcd}[column sep=small, row sep=small, nodes={font=\scriptsize}]
U_{03} \arrow[r, "{\{1, \omega\}}"] \arrow[dr, "\mathbb{D}'_3"'] & (U_{0,3}, \{1, \omega\}) \arrow[d, "\mathbb{D}_3"'] \arrow[dr, "\mathrm{gerbe}"] & \\
 & (U_{0,3}, \mathbb{D}'_3) \arrow[d, "\mathrm{gerbe}"'] & (U_{0,3}) \arrow[dl, "\mathbb{D}_3"] \\
 & (U_{0,3}, \mathbb{D}_3) &
\end{tikzcd}

\note{redessiné sur une grille. Sur la page, la flèche marquée
$\mathbb{D}_3$ descend en oblique de $(U_{0,3}, \{1, \omega\})$ vers
$(U_{0,3}, \mathbb{D}'_3)$, et la flèche « gerbe » descend verticalement de
$(U_{0,3}, \{1, \omega\})$ vers $(U_{0,3})$}

\begin{tikzcd}[column sep=small, row sep=small, nodes={font=\scriptsize}]
(\mathcal{D}, \Pi_{0,3}) \arrow[r, Rightarrow, "\mathrm{gal}"] \arrow[dr, "\mathrm{gal}"'] & (\mathcal{D}, \Pi_{0,3} \times \{1, \omega\}) \arrow[d, "\mathrm{gal}"'] \arrow[dr, "2\text{-}\mathrm{gerbe}"] & \\
 & (\mathcal{D}, \mathrm{SL}(2, \mathbb{Z})) \arrow[d, "2\text{-}\mathrm{gerbe}"'] & (\mathcal{D}, \Pi_{0,3}) \arrow[dl, "\mathrm{gal}"] \\
 & (\mathcal{D}, \mathrm{SL}(2, \mathbb{Z})') &
\end{tikzcd}

\note{même disposition que le diagramme précédent, dont celui-ci est le
calque. Sous $(\mathcal{D}, \Pi_{0,3})$ : « $\simeq \mathcal{D}/\Pi_{0,3}$ ».
La flèche horizontale porte « gal » au-dessus et « $\{1, \omega\}$ »
au-dessous. Les étiquettes des trois flèches « gal » obliques et verticale
nomment le groupe : $\Pi^{\mathbb{D}'}_{0,3}$ ou $\mathbb{D}'_3$
\uncertain{} (surchargé) pour la flèche de gauche, $\mathbb{D}_3$ pour celle
du milieu (sur un « $\{1, \omega\}$ » et un $\Pi$ biffés), $\mathbb{D}_3$
pour celle de droite. Le $2$ de « 2-gerbe » est surchargé ; « gerbe » est
souligné sur la flèche de droite. Le $\mathcal{D}$ est une capitale
cursive, lue ici comme le demi-plan de Poincaré \uncertain{}}

\section{Extensions de $\mathbb{D}_3$ et de $\Pi^{D}_{0,3}$}

\note{titre de l'édition. Page 28, au crayon : un grand dessin en
perspective, deux cubes superposés, que traversent trois longs traits
obliques recourbés en crosse, sans rien biffer. Il est redessiné ici sur une
grille ; les positions relatives sont celles du dessin autant qu'elles se
laissent rendre, et la place de quelques flèches est incertaine}

\page{28}

\begin{tikzcd}[column sep=tiny, row sep=small, nodes={font=\scriptsize}]
 & \mathrm{SL}(2, \mathbb{Z}) \arrow[r] \arrow[d] & \Pi^{D}_{0,3} \arrow[d] & & \\
\widetilde{E} \arrow[r] \arrow[d] \arrow[dr, dashed] & \mathbb{D}_3 \times_{\mathbb{Z}/12\mathbb{Z}} \mathbb{Z}/12\mathbb{Z} \simeq \mathbb{D}'_3 \times_{\mathbb{Z}/12\mathbb{Z}} \mathbb{Z}/6\mathbb{Z} \arrow[r] & \mathbb{D}_3 \times_{\mathbb{Z}/12\mathbb{Z}} \mathbb{Z}/6\mathbb{Z} = E & & \\
\widetilde{\mathbb{D}}_3 \arrow[dr, Rightarrow] & \mathbb{Z} \arrow[r, dashed] & \mathbb{Z}/12\mathbb{Z} \arrow[r, dashed, "\mathrm{can}"] \arrow[d, dashed] & \mathbb{Z}/6\mathbb{Z} \arrow[d, "\mathrm{can}"] & \\
 & \mathbb{D}'_3 \arrow[r, Rightarrow] \arrow[d] & \mathbb{D}_3 \arrow[r, "\mathrm{can}"] & \mathbb{Z}/2\mathbb{Z} & \\
\mathbb{Z} \arrow[r] & \mathbb{Z}/4\mathbb{Z} \arrow[urr, "\mathrm{can}"'] & & &
\end{tikzcd}

\note{sur la page les nœuds sont disposés en perspective, non sur une
grille. Lecture des flèches : $\widetilde{E} \to \mathbb{D}_3
\times_{\mathbb{Z}/12\mathbb{Z}} \ldots$ (horizontale) ; de
$\widetilde{E}$ partent aussi une flèche vers $\widetilde{\mathbb{D}}_3$ et
une flèche en tirets vers le premier $\mathbb{Z}$ ; une autre flèche en
tirets va de $\widetilde{\mathbb{D}}_3$ \uncertain{} au second
$\mathbb{Z}$ ; $\widetilde{\mathbb{D}}_3 \Rightarrow \mathbb{D}'_3$ en trait
double et gras ; $\mathbb{D}'_3 \to \mathbb{D}_3$ en trait double ;
$\mathbb{Z} \cdots \to \mathbb{Z}/12\mathbb{Z}$ en pointillés,
$\mathbb{Z}/12\mathbb{Z} \dashrightarrow \mathbb{Z}/6\mathbb{Z}$ marquée
« can » ; $\mathbb{Z}/12\mathbb{Z} \dashrightarrow \mathbb{Z}/4\mathbb{Z}$ ;
$\mathbb{D}'_3 \to \mathbb{Z}/4\mathbb{Z}$ ; $\mathbb{D}_3$,
$\mathbb{Z}/6\mathbb{Z}$ et $\mathbb{Z}/4\mathbb{Z}$ vont vers
$\mathbb{Z}/2\mathbb{Z}$, chaque flèche marquée « can » ; le produit fibré
du milieu descend vers $\mathbb{D}'_3$ et vers $\mathbb{Z}/12\mathbb{Z}$,
et $E$ descend vers $\mathbb{D}_3$ et vers $\mathbb{Z}/6\mathbb{Z}$. Une
flèche horizontale part du bord gauche de la page, en haut, vers
$\mathrm{SL}(2, \mathbb{Z})$, avec un petit trait vertical. Les indices des
deux premiers produits fibrés se lisent $\mathbb{Z}/12\mathbb{Z}$ ; le
facteur de droite du premier est lu $\mathbb{Z}/12\mathbb{Z}$
\uncertain{}}

\begin{tikzcd}[column sep=small, row sep=small, nodes={font=\scriptsize}]
\mathbb{Z}(\frac{\omega}{2}) \arrow[r] \arrow[d, Rightarrow] & \widetilde{\Pi}^{\circ}_{0,3} \arrow[r, Rightarrow] \arrow[d] & \Pi_{0,3} \arrow[d, hook] & \\
\mathbb{Z}(\omega) \arrow[r] \arrow[d, no head, "\parallel" description] & \widetilde{\Pi}'_{0,3} \arrow[r] \arrow[d] & \Pi'_{0,3} = \Pi_{0,3} \times \{1, \dot\omega\} \arrow[r] \arrow[d] & \Pi_{03} \arrow[d, hook] \\
\mathbb{Z} \arrow[r] \arrow[d, no head, "\parallel" description] & \widetilde{\mathrm{SL}}(2, \mathbb{Z}) = \widetilde{\Pi}^{D}_{0,3} \arrow[r] \arrow[d] & \mathrm{SL}(2, \mathbb{Z}) \arrow[r] \arrow[d] & \Pi^{D}_{0,3} \arrow[d] \\
\mathbb{Z} \arrow[r] \arrow[d, no head, "\parallel" description] & \widetilde{\mathbb{D}}_3 \arrow[r] \arrow[d] & \mathbb{D}'''_3 \arrow[r, "2"] \arrow[d, "3"] & \mathbb{D}''_3 \arrow[d, "3"] \\
\mathbb{Z} \arrow[r, "12"] & \mathbb{Z} \arrow[r] & \mathbb{Z}/12\mathbb{Z} \arrow[r, "2"] & \mathbb{Z}/6\mathbb{Z}
\end{tikzcd}

\begin{tikzcd}[column sep=small, row sep=small, nodes={font=\scriptsize}]
\mathbb{D}'''_3 \arrow[r, "2"] \arrow[d, "3"'] \arrow[dr, "3"] & \mathbb{D}''_3 \arrow[d, "3"] \arrow[dr, "3"] & \\
\mathbb{Z}/12\mathbb{Z} \arrow[dr, "3"'] & \mathbb{D}'_3 \arrow[r, "2"] \arrow[d, "3"] & \mathbb{D}_3 \arrow[d, "3"] \\
 & \mathbb{Z}/4\mathbb{Z} \arrow[r, "2"] & \mathbb{Z}/2\mathbb{Z}
\end{tikzcd}

\note{les deux diagrammes précédents n'en font qu'un sur la page : le
second est la face avant du cube dont $\mathbb{D}'''_3$,
$\mathbb{D}''_3$, $\mathbb{Z}/12\mathbb{Z}$, $\mathbb{Z}/6\mathbb{Z}$ sont
la face arrière ; il est détaché ici pour la lisibilité, et la flèche
$\mathbb{Z}/12\mathbb{Z} \to \mathbb{Z}/6\mathbb{Z}$ (marquée $2$) et
$\mathbb{Z}/6\mathbb{Z} \to \mathbb{Z}/2\mathbb{Z}$ (marquée $3$)
appartiennent aux deux. En outre : $\mathrm{SL}(2, \mathbb{Z})$ envoie une
flèche oblique vers $\mathbb{D}'_3$, $\Pi^{D}_{0,3}$ une flèche oblique vers
$\mathbb{D}_3$ ; de $\Pi'_{0,3}$ part une flèche vers $\Pi_{03}$ et, depuis
$\Pi_{03}$, une flèche courbe grasse vers $\Pi^{D}_{0,3}$. Les chiffres
$2$ et $3$ portés par les flèches sont de sa main ; sous $\mathbb{D}''_3$, un
« $18$ » \uncertain{} très pâle. L'étiquette de la flèche
$\mathbb{Z}(\omega/2) \Rightarrow \mathbb{Z}(\omega)$ est entourée ; le
$\widetilde{\Pi}^{\circ}_{0,3}$ est suivi d'un second $\widetilde{\Pi}$
effacé. Le mot devant le $\mathbb{Z}$ du bas à gauche est illisible}



\note{bas de la page 28 : trois suites exactes, séparées de la dernière par
un trait}

\[
\begin{array}{ccccccccc}
1 & \to & \mathbb{Z} & \to & \widetilde{\Pi}^{D}_{0,3} & \to & \Pi^{D}_{0,3} & \to & 1 \\
 & & \downarrow{\scriptstyle 4} & & \downarrow & & \downarrow & & \\
1 & \to & \mathbb{Z} & \to & \widetilde{\Pi}^{D}_{0,3}(4) & \to & \Pi^{D}_{0,3} & \to & 1
\end{array}
\]

\note{l'étiquette $4$ de la première flèche verticale est récrite sur un
autre signe ; la dernière flèche de la deuxième ligne est tracée double.
Sous le trait, à l'aplomb : « $\widetilde{\Pi}^{D}_{0,3}$ \quad
$\Pi^{D}_{0,3} \to 1$ »}

\[
\begin{array}{ccccccccc}
 & & & & & & \Pi^{D}_{0,3} & & \\
 & & & & & & \uparrow & & \\
1 & \to & \overset{\omega}{\mathbb{Z}} & \to & \widetilde{\Pi}^{D}_{0,3}(4) & \to & \Pi^{D}_{0,3} \times \mathbb{Z}/4\mathbb{Z} & \to & 1 \\
 & & & & & & \uparrow & & \\
1 & \to & \mathbb{Z} & \to & ? & \to & \struck{\ill{}}\ \mathrm{SL}(2, \mathbb{Z}) & \to & 1
\end{array}
\]

\note{le point d'interrogation tient lieu du terme médian de la dernière
suite ; il est de sa main}

\section{Sections de $\mathcal{O}(1)$ et changement d'orientation}

\note{titre de l'édition. Page 29, encre noire ; un long trait de crayon la
traverse en oblique. La moitié supérieure est à l'endroit ; la moitié
inférieure est écrite tête-bêche et se lit feuille retournée ; une colonne
est écrite en travers dans la marge de gauche}

\page{29}

$V_J$ \struck{\ill{}} $= \Gamma(X, \underline{\mathcal{O}}(1))$
\uncertain{engendré} par les $e_i - e_j$

\struck{$P(e_i - e_j)$} \ill{} \ill{} \ill{} \uncertain{ces}
\uncertain{sections} \ill{} \ill{} \uncertain{elles}

\uncertain{permutations} \uncertain{entre} \uncertain{elles}

\note{à gauche de ces lignes, un cercle portant trois petits traits
(repères), avec un point marqué $e$ à gauche. La flèche qui part de
$\Gamma(X, \underline{\mathcal{O}}(1))$ ramène à la première ligne}

\[
\struck{e_i - e_j \quad \sigma(e)} \qquad e_i - e_j =
\qquad\qquad
\underline{\mathcal{O}}_X(U) \ni \xi_{ij}(x)
\]
\[
\xi_i \qquad\qquad
\lambda \qquad \lambda^6 = \prod_{i,j} \xi_{ij}(x)
\]

\note{sous $\underline{\mathcal{O}}_X(U)$, un signe $\Psi$ \uncertain{}
au-dessus de $\lambda$ ; les indices de $\prod$ sont en pointillé}

\[
(e_i - e_j)(\underline{Q_k}) \qquad \struck{e_i - e_j}
\]
\[
\mathbb{Z} \to \widetilde{\Pi}_{0,3} \to \Pi^{\struck{\ill{}}}_{03} \to 1
\]
\[
V(\omega) \overset{!!}{\simeq} V \qquad \varphi \in V \qquad
e_i - e_j(\struck{\ill{}}_K
\]
\[
e_i - e_j \bmod L_{Q_K}
\qquad
\begin{array}{ccc}
\frac12 & \frac12 & 1 \\[2pt]
-\frac12 & -\frac12 & -1
\end{array}
\qquad \Bigl(\frac12\Bigr)^4 \qquad \sqrt[3]{1/2}
\]

\note{les nombres sont dispersés à droite de la ligne, sans relation écrite
entre eux}

\note{ce qui suit est écrit tête-bêche (moitié inférieure de la page,
feuille retournée)}

On a vu qu'une orientation d'une \struck{\uncertain{complexe}}
\ill{} combinatoire $A$ donne lieu : aux opérations $\vec\rho_s$,
$\vec\rho_f$ sur $\bar A$, satisfaisant
\[
\vec\rho_s^{\,-1}\,\vec\rho_f = \sigma .
\]
Si on prend l'orientation inverse, quelles sont les opérations
correspondantes $\vec\rho\,'_s$, $\vec\rho\,'_f$ ? Si $\bar a \in \bar A$
provient de $a \in A^L_+$, il provient de $\sigma_n a \in A^L_-$, et on a
\[
\rho_s(\sigma_n a) = \sigma_n(\sigma_n \underbrace{\rho_s}_{\sigma_n \varepsilon = \rho_s^{-1}} \sigma_n\, a)
= \sigma_n(\rho_s^{-1} a) .
\]
donc la classe de $\bar A \simeq A^L/\sigma_n$ est celle de
$\rho_s^{-1} a$, i.e. $\vec\rho_s^{\,-1} \bar a$. \struck{\ill{}} \ill{}
…
\[
\struck{\rho_f(\sigma_n a) = \sigma_n(\sigma_n \rho_f \sigma_n\, a)}
\]
\[
\vec\rho\,'_s = \vec\rho_f^{\,-1}
\]
d'où
\[
\vec\rho\,'_f = \vec\rho\,'_s\,\sigma = \vec\rho_s^{\,-1}\sigma
= \sigma\,\vec\rho_f^{\,-1}\,\sigma = \mathrm{int}(\sigma)\,\vec\rho_f^{\,-1}
\]

\note{au-dessus du mot biffé de la première ligne, un mot inséré,
illisible. L'accolade sous $\rho_s$ porte « $\sigma_n \varepsilon =
\rho_s^{-1}$ », l'indice de $\rho$ se lisant $s$ ou $f$. Sous la ligne
biffée, une accolade marquée « $\sigma_n \varepsilon \sigma$ ». Dans
« $\vec\rho\,'_s = \vec\rho_f^{\,-1}$ » l'indice de droite est lu $f$ :
il contredit la ligne précédente ($\vec\rho_s^{\,-1}$), et la suite
utilise $\vec\rho_s^{\,-1}$}

\marginal{$\tilde R_K - \tilde R_j$ ? \quad \ill{} \quad
$e_j - e_i - (e_i - e_k)$}

\note{colonne écrite en travers dans la marge de gauche, lisible la feuille
tournée d'un quart de tour}

\section{Les extensions $\widetilde\Pi$, $\Pi'$ et $\mathbb{D}'_3$}

\note{titre de l'édition. Page 30, encre noire ; les $\mathbb{D}$ y sont
écrits comme des $D$ ordinaires}

\page{30}

\[
1 \hookrightarrow \mathbb{Z} \to \widetilde\Pi \to \Pi \to 1
\]

\note{la suite est soulignée d'un trait qui se recourbe au bout. De $\Pi$
descend une flèche vers $\mathbb{D}''_3$, écrit sur un
« $\mathbb{Z}/6$ » \uncertain{} biffé. Au-dessous : « $\Pi_0 = \Pi_0$ »,
d'où partent deux flèches, vers $\Pi'$ et vers $\Pi$}

\begin{tikzcd}[column sep=small, row sep=small, nodes={font=\scriptsize}]
\mathbb{Z} \arrow[r] \arrow[d, no head, "\parallel" description] & \widetilde\Pi \arrow[r] \arrow[d] & \Pi' \arrow[r] \arrow[d] \arrow[ddr, bend left=20] & \Pi \arrow[d] \arrow[dr] & \\
\mathbb{Z} \arrow[r] \arrow[d, no head, "\parallel" description] & \widetilde{\mathbb{D}}_3 \arrow[r] \arrow[d] & \mathbb{D}'''_3 \arrow[r] \arrow[d] \arrow[dr] & \mathbb{D}''_3 \arrow[d] \arrow[r, Rightarrow] & \mathbb{D}_3 \arrow[dd] \\
\mathbb{Z} \arrow[r, "12"] & \mathbb{Z} \arrow[r] & \mathbb{Z}/12\mathbb{Z} \arrow[d] & \mathbb{D}'_3 \arrow[ur] \arrow[d, no head] & \\
 & & \mathbb{Z}/4\mathbb{Z} \arrow[rr] & \mathbb{Z}/6\mathbb{Z} \arrow[r] & \mathbb{Z}/2\mathbb{Z}
\end{tikzcd}

\note{redessiné sur une grille : sur la page, $\mathbb{D}'_3$ est placé
entre $\mathbb{D}'''_3$ et $\mathbb{Z}/6\mathbb{Z}$, et sa flèche vers
$\mathbb{D}_3$ passe au-dessus de $\mathbb{Z}/6\mathbb{Z}$ ; le trait
vertical qui descend de $\mathbb{D}'_3$ coupe la flèche
$\mathbb{Z}/12\mathbb{Z} \to$ et va jusqu'à $\mathbb{Z}/4\mathbb{Z}$. La
flèche $\mathbb{D}''_3 \Rightarrow \mathbb{D}_3$ est tracée double ;
$\mathbb{Z}/12\mathbb{Z} \to \mathbb{Z}/4\mathbb{Z}$ est un trait
oblique sans tête}

\[
\begin{array}{ccccccccc}
 & & & & \Pi_0 & & & & \\
 & & & & \downarrow & & & & \\
1 & \to & \mathbb{Z} & \to & \widetilde\Pi(4) & \to & \Pi \times \mathbb{Z}/4\mathbb{Z} & \to & 1
\end{array}
\]

\note{à gauche de la flèche verticale, « $\struck{\widetilde\Pi}$
$\struck{\mathbb{Z}} \to$ », biffés. Sous $\Pi \times \mathbb{Z}/4\mathbb{Z}$,
une flèche vers un $\struck{\mathbb{D}_3 \times \mathbb{Z}/4\mathbb{Z}}$
raturé en boucles, et une flèche oblique vers
$\mathbb{D}_3 \times \mathbb{Z}/4\mathbb{Z} \supset \mathbb{D}'_3$ ; un
long trait courbe descend de la rature vers le bas}

\[
\begin{array}{ccccccccc}
1 & \to & \Pi_0 & \to & \Pi \times \mathbb{Z}/4\mathbb{Z} & \to & \mathbb{D}_3 \times \mathbb{Z}/4\mathbb{Z} & \to & 1 \\
 & & \parallel & & \uparrow & & \updownarrow & & \\
1 & \to & \Pi_0 & \to & \Pi' & \longrightarrow & \mathbb{D}'_3 & \to & 1
\end{array}
\]

\[
\Pi' \subset \Pi \qquad\qquad
\struck{\widetilde\Pi \times \mathbb{Z}/4} \quad
\widetilde\Pi(4) \to \Pi \times \mathbb{Z}/4\mathbb{Z}
\]

\note{la flèche $\Pi' \to \Pi \times \mathbb{Z}/4\mathbb{Z}$ est tracée avec
une boucle en bas (inclusion) ; celle de droite a une tête à chaque bout}

\section{Orbites de $\rho$ et quadruples $(A, \sigma, A', \rho)$}

\note{titre de l'édition. Page 31, à l'encre bleue, rédaction rapide ; un
long trait de crayon la traverse en oblique sans rien biffer. Beaucoup de
mots de liaison ne se lisent pas}

\page{31}

\emph{Cas a)} La construction \uncertain{ne} \ill{} se continue
indéfiniment, i.e. on \uncertain{finit} par avoir $\sigma a_n \notin A'$,
on \uncertain{pose} alors $\sigma a_n = \rho\, a_0$.

\note{au-dessus de « construction », un mot inséré, lu \uncertain{ne}}

\emph{Cas b)} La construction se prolonge indéfiniment.
(\struck{\ill{} par $a_i$ \ill{} un cycle fini} \uncertain{donc}
\uncertain{un} cycle infini). On ne définit pas $\rho\, \bar a$ alors.

* On supposera désormais que $\rho$ est défini
\uncertain{\struck{partout}} sur tout $A$, i.e. \ill{} que
l'orientation inverse (ce sera le cas de $(A, \sigma, A', \rho)$). Alors
$\rho$ sur $A$ est un \emph{automorphisme}. Soit
$\bar S = A/\rho \supset S$, $\hat S \subset \bar S$ la partie de $\bar S$
correspondant aux orbites finies. On va compactifier partiellement $X$ en
$\hat X$ \struck{\ill{}} \uncertain{ajoutant} les \uncertain{points} de
$\hat S - S$ (ainsi \ill{} qui : $K^{\circ}$ \ill{}, l'image inverse
$\hat A'$ de $\hat S$ dans $A$ (\ill{} \ill{} \ill{} $\rho$ et $\sigma$
\ill{}) \struck{\ill{} \ill{} \ill{} : $\hat A$ \ill{}} \ill{} $\hat A'$.

\note{les mots insérés au-dessus des lignes (« $\hat A'$ », « $K_{0,3}$ »
\uncertain{}, « \uncertain{contient} $A'$ et ») sont reliés par des traits
à leur place ; la phrase ne se laisse pas reconstruire}

\marginal{* \ill{} de circuits définis dans un seul \uncertain{sens}}

\marginal{\uncertain{stabilité} de l'\uncertain{axiome} $A' = \hat A'$ ?
\ill{} mais \struck{\ill{}}}

\note{la première note est écrite de biais dans la marge de gauche, à la
hauteur de l'astérisque ; la seconde est entourée d'une ligne fermée,
en bas à gauche, et reliée par deux traits au texte}

Les quadruples $(A, \sigma, A', \rho)$ satisfaisant aux axiomes
[\uncertain{et} tels que $A' = \hat A'$], dits \ill{},
\uncertain{s'identifient} aux \struck{couples} triples $(A, \sigma, \rho)$
d'un ens. \emph{\uncertain{fini}} avec un automorphisme $\sigma$ d'ordre
$2$ \emph{sans pt fixe}, et un automorphisme $\rho$
\emph{\uncertain{quelc.}}. On retrouve $A'$ comme l'une des \ill{} de
$A$ d'orbite finie sous $\rho$ --- mais --- telle que
\[
A' \cup \sigma A' = A .
\]

\note{« et tels que $A' = \hat A'$ » est inséré au-dessus de la ligne, sous
une accolade}

\section{Classes d'extensions de $\mathbb{Z}/p \times \mathbb{Z}/q$ par $\mathbb{Z}$}

\note{titre de l'édition. Page 32 à l'encre brune}

\page{32}

\begin{tikzcd}[column sep=small, row sep=small, nodes={font=\scriptsize}]
0 \arrow[r] & \overset{\omega}{\mathbb{Z}} \arrow[r, "p"] \arrow[d, "q"] & \overset{\widetilde\Pi_0}{\mathbb{Z}} \arrow[r, "\alpha\varphi_p"] \arrow[d] & \overset{\Pi_0}{\mathbb{Z}/p\mathbb{Z}} \arrow[r] \arrow[d, no head, "\parallel" description] & 0 \\
0 \arrow[r] & \mathbb{Z} \arrow[r] & \widetilde\Pi_0(q) \arrow[r] \arrow[d] & \mathbb{Z}/p\mathbb{Z} \arrow[r] & 0 \\
 & & \mathbb{Z}/q\mathbb{Z} & &
\end{tikzcd}

\note{les $\widetilde\Pi_0$, $\Pi_0$ sont écrits au-dessus des groupes,
reliés par un signe $=$ vertical. La flèche marquée $\alpha\varphi_p$ est
grasse. Sous le $\mathbb{Z}$ de la deuxième ligne, un « $\varpi$ ». La
flèche $\widetilde\Pi_0(q) \to \mathbb{Z}/q\mathbb{Z}$ porte : « quotient
par $\widetilde\Pi_0$, \uncertain{can.} par $\mathbb{Z}\varpi$ »}

\[
\varphi_p : \mathbb{Z} \to \mathbb{Z}/p\mathbb{Z} \quad \text{homom. can.}
\qquad
\alpha \in (\mathbb{Z}/p\mathbb{Z})^{*}
\qquad
\struck{p \in (\mathbb{Z}/q\mathbb{Z})^{*}}
\]

\begin{tikzcd}[column sep=small, row sep=small, nodes={font=\scriptsize}]
 & \overset{\omega}{\mathbb{Z}} \arrow[r] \arrow[dl, Rightarrow] & \widetilde\Pi_0 \arrow[r] \arrow[dl] & (\mathbb{Z}/p\mathbb{Z} \times 0) \arrow[dl] & \\
0 \arrow[r] & \overset{\omega}{\mathbb{Z}} \arrow[r] & \widetilde\Pi_0(q) \arrow[r] & \mathbb{Z}/p\mathbb{Z} \times \mathbb{Z}/q\mathbb{Z} \arrow[r] & 0 \\
 & \underset{\omega}{\mathbb{Z}} \arrow[r] \arrow[u, Rightarrow, "q"] & \underset{\varpi}{\mathbb{Z}} \arrow[r] \arrow[u] & 0 \times \mathbb{Z}/q\mathbb{Z} \arrow[u] &
\end{tikzcd}

\note{sur la page les trois lignes sont décalées en perspective et les
flèches entre elles sont obliques. À droite de la première ligne :
« $\alpha\eta_p$ ». La flèche $\mathbb{Z} \to \mathbb{Z}$ du bas porte un
signe surchargé \ill{}}

\[
\begin{array}{ccccccccc}
0 & \to & \mathbb{Z} & \to & \widetilde\Pi_0(q) & \to & \mathbb{Z}/p\mathbb{Z} \times \mathbb{Z}/q\mathbb{Z} & \to & 0 \\
 & & \parallel & & \updownarrow & & \uparrow{\scriptstyle\delta} & & \\
0 & \to & \mathbb{Z} & \to & \widetilde\Pi'_0 & \longrightarrow & \mathbb{Z}/\mu\mathbb{Z} & \to & 0
\end{array}
\]

\struck{d'extension \ill{}}
\[
\xi' = \mathrm{cl}(\widetilde\Pi'_0) \in
\mathrm{Ext}^1(\mathbb{Z}/\mu\mathbb{Z}, \mathbb{Z}) \simeq \mathbb{Z}/\mu\mathbb{Z}
\]
\[
\xi' = \delta^{*}\bigl(\underbrace{\mathrm{cl}(\widetilde\Pi_0(q))}_{\xi \in
\mathrm{Ext}^1(\mathbb{Z}/p\mathbb{Z} \times \mathbb{Z}/q\mathbb{Z},\, \mathbb{Z})}\bigr)
\]
\[
\xi = \mathrm{pr}_1^{*}\bigl(\underbrace{\mathrm{cl}(\widetilde\Pi_0)}_{\alpha\eta_p}\bigr)
+ \mathrm{pr}_2^{*}\bigl(\underbrace{\mathrm{cl}(\mathbb{Z})}_{\eta_q}\bigr)
\]

\note{le symbole à gauche du premier $=$ est récrit sur une rature ;
l'argument de $\mathrm{pr}_2^{*}$ est inachevé, une parenthèse ouverte
avant un blanc}

\[
\delta_p : \mathbb{Z}/\mu\mathbb{Z} \xrightarrow{\ \mathrm{can}\ } \mathbb{Z}/p\mathbb{Z}
\qquad
\delta_q : \mathbb{Z}/\mu\mathbb{Z} \xrightarrow{\ \mathrm{can}\ } \mathbb{Z}/q\mathbb{Z}
\qquad
\xi' \struck{= q'\eta_\mu + p'\eta_\mu}
\]
\[
\delta_p^{*}(\eta_p) = q'\eta_\mu
\qquad
\delta_q^{*}(\eta_q) = p'\eta_\mu
\qquad
\delta_p^{*}(\alpha\eta_p) = q'\alpha\,\eta_\mu
\]

\section{$S\mathcal{T}_{11}$ et $S\Pi^{D}_{0,3}$}

\note{titre de l'édition. Page 33, encre brune. Le haut est à l'endroit ;
le reste de la page est écrit tête-bêche}

\page{33}

\struck{4) $\widetilde{\mathrm{SL}}(2, \mathbb{Z}) \simeq \Pi_{\ill{}}$
$S\mathcal{T}_{\ill{}}$}

\uncertain{1)}
\[
S\mathcal{T}_{11} \simeq \widetilde{\mathrm{SL}}(2, \mathbb{Z})
\simeq \{\rho, \sigma \mid \underbrace{\rho^3 = \sigma^2}_{\omega}\}
\simeq \underset{\{\rho, \sigma \mid \rho^3 = \sigma^2 = 1\}}{\underset{\parallel}{\Pi^{D}_{03}}}
\times_{\mathbb{Z}/6\mathbb{Z}} \mathbb{Z}
\]
\[
\simeq S\Pi^{D}_{0,3}
\]

\note{l'indice du produit fibré est écrit sur un $\mathbb{Z}$ biffé. Un
mot biffé est relié par une courbe à « $\rho \mapsto 2$,
$\sigma \mapsto 3$ », écrit sous le produit fibré}

\note{la suite de la page est écrite tête-bêche ; elle est donnée
feuille retournée}

\[
\begin{array}{ccccccccc}
 & & \mathbb{Z} & \to & \widetilde\Pi(4) & \to & \Pi \times \mathbb{Z}/4\mathbb{Z} & \to & 1 \\
 & & & & \uparrow{\scriptstyle ?} & & \downarrow & & \\
 & & & & \widetilde\Pi & \longrightarrow & \Pi' & &
\end{array}
\]

\note{la flèche verticale de gauche porte « ? » et un mot souligné
\ill{}. Au-dessous, une courbe entoure « $\rho \mapsto 2$,
$\sigma \mapsto 3$ » (les mêmes mots que plus haut, lus dans l'autre
sens) ; à droite, « $\Pi_0 \times \mathbb{Z}$ » \uncertain{}}

\begin{tikzcd}[column sep=small, row sep=small, nodes={font=\scriptsize}]
\widetilde\Pi \arrow[r, no head] \arrow[d, no head] & \Pi \arrow[dd, "\mathbb{D}''_3" description] \arrow[r, no head] & \Pi' \arrow[r] \arrow[d, Rightarrow] & \Pi \arrow[d] \\
\mathbb{Z} \arrow[rr] & & \mathbb{Z}/12\mathbb{Z} \arrow[r] \arrow[d] & \mathbb{Z}/6\mathbb{Z} \arrow[d] \\
 & \mathbb{Z}/4\mathbb{Z} & \mathbb{Z}/4\mathbb{Z} \arrow[r] & \mathbb{Z}/2\mathbb{Z}
\end{tikzcd}

\note{redessiné sur une grille. Le trait vertical qui descend du $\Pi$
de la première ligne, marqué $\mathbb{D}''_3$ à mi-hauteur, jusqu'à un
« $\mathbb{Z}/4$ » \uncertain{} coupe les deux flèches horizontales ; la
flèche $\Pi' \Rightarrow \mathbb{Z}/12\mathbb{Z}$ est tracée double et
grasse. À droite : « $\Pi_0 = \mathbb{Z}/6\mathbb{Z}$ »}

\note{la page porte encore, tête-bêche et d'une autre main, à l'encre
rouge, un bref mot adressé à lui et daté 27-11-81 ; il n'est pas
transcrit}

\section{Le caractère $\Pi^{D}_{0,3} \to \mathbb{Z}/6\mathbb{Z}$}

\note{titre de l'édition. Page 34, sur papier quadrillé, encre noire}

\page{34}

\[
(p' + q'\alpha)\,\eta_\mu
\]
\[
p = 6, \quad q = 4 \qquad (\alpha = -1\ ?)
\]
\[
\struck{p'}\ \mu = 12 \qquad p' = 3, \quad q' = 2
\]
\[
\struck{p'+}\ (3 - 2)\,\eta_\mu \quad !
\]

\note{les valeurs $p' = 3$, $q' = 2$ sont celles de $\mu/q$ et $\mu/p$}

\[
\begin{array}{rcl}
\Pi^{D}_{0,3} & \xrightarrow{\ -\alpha'\ } & \mathbb{Z}/6\mathbb{Z} \\
\rho & \longmapsto & -2 \bmod 6 \\
\sigma & \longmapsto & 3 \bmod 6 \\
\varepsilon_0 = \sigma\rho & \longmapsto & 1 \bmod 6 \\
\rho_0 = \varepsilon_0^{\,2} & \longmapsto & 2 \bmod 6
\end{array}
\]

\note{l'étiquette « $-\alpha'$ » est entourée, écrite sur un « $q$ »
\uncertain{}}

défini directement par l'action de $\Pi^{D}_{0,3}$ sur le
\struck{revêtement} \uncertain{revêtement} principal de
$(U_{03}, \mathfrak{S}_3)$ de groupe $\mathbb{Z}/6\mathbb{Z}$, défini par
la multiplication \struck{\ill{}} \uncertain{canonique}
\struck{\ill{}} $\mathcal{O}_{U_{0,3}}(1)^{*} = PU_{03}$

\note{« canonique » est inséré au-dessus des mots biffés. Entre les deux
colonnes, quelques lettres au crayon très pâle (« $T$ »,
« $\varepsilon'$ », « $t'h$ » \uncertain{}), sans lien visible}

\[
\widetilde\Pi^{D}_{03}
\]

\[
\struck{\ill{}}\quad U'_{03} \subset PU_{03}
\qquad
\mu_6\text{-}\uncertain{\text{revêt.}}\ \struck{\ill{}}\,/U_{03}
\]

\begin{tikzcd}[column sep=small, row sep=small, nodes={font=\scriptsize}]
\Pi^{D}_{0,3} \arrow[r, "\alpha'"] & \mathbb{Z}/6\mathbb{Z} \\
\pi_1(PU_{03}, \mathfrak{S}_3) = \widetilde\Pi^{D}_{03} \arrow[u] \arrow[r] & \mathbb{Z} \arrow[u] \\
\varpi' \in \mathbb{Z} \arrow[u] \arrow[r, "="] & \mathbb{Z} \arrow[u, "6"']
\end{tikzcd}

\note{l'étiquette $\alpha'$ est entourée}

\[
\begin{array}{ccccccccc}
1 & \to & \mathbb{Z} & \to & \widetilde\Pi^{D}_{03}(4) & \to & \Pi^{D}_{0,3} \times \mathbb{Z}/4\mathbb{Z} & \to & 1 \\
 & & \parallel & & \parallel & & \updownarrow & & \\
1 & \to & \mathbb{Z} & \to & \widetilde{\mathcal{T}}_{11} & \Longrightarrow & \mathcal{T}_{1,1} = \Pi^{D}_{0,3} \times_{\mathbb{Z}/12\mathbb{Z}} \mathbb{Z}/4\mathbb{Z} & &
\end{array}
\]

\note{le signe $=$ vertical sous $\widetilde\Pi^{D}_{03}(4)$ porte, entre
parenthèses ouvertes : « $\pi_1(PU_{03}, \mathfrak{S}_3 \times \mu_4)$ »}

\[
\struck{\ill{}}\qquad
\widetilde\Pi^{D}_{0,3} \ni
\begin{cases}
\tilde\rho = (\rho, -2) \\
\tilde\sigma = (\sigma, -3)
\end{cases}
\qquad
\tilde\rho^{\,3} = \tilde\sigma^{2} = \varpi'
\]

\note{les exposants de $\tilde\rho$ et $\tilde\sigma$, dans les deux
lignes, sont écrits sur d'autres chiffres raturés, illisibles}

\[
\underset{\substack{\downarrow \\ \rho}}{\tilde\rho},\ \underset{\substack{\downarrow \\ \sigma}}{\tilde\sigma}
\in \widetilde\Pi^{D}_{0,3}
= \{\tilde\rho, \tilde\sigma \mid \underbrace{\tilde\rho^{\,3} = \tilde\sigma^{2}}_{\varpi' = \omega}\ \struck{\ill{}}\}
\simeq \widetilde{\mathcal{T}}_{11}
\]

\page{35}

\begin{tikzcd}[column sep=small, row sep=small, nodes={font=\scriptsize}]
1 \arrow[r] & \overset{\omega' = \varpi^{-1}}{\mathbb{Z}} \arrow[r] \arrow[d, no head, "\parallel" description] & S\Pi^{D}_{0,3} \arrow[r] \arrow[d] & \Pi^{D}_{03} \arrow[r] \arrow[d, hook] & 1 \\
1 \arrow[r] & \mathbb{Z} \arrow[r] & S\Pi^{D}_{03}(4) \arrow[r] & \Pi^{D}_{03} \times \mathbb{Z}/4\mathbb{Z} \arrow[r] & 1 \\
1 \arrow[r] & \underset{\omega^2}{\mathbb{Z}} \arrow[r] \arrow[u] & S\mathcal{T}_{11} \arrow[r] \arrow[u] & \mathcal{T}_{11} \arrow[u, leftrightarrow] &
\end{tikzcd}

\note{au-dessus de $S\Pi^{D}_{0,3}$ : « $\{\rho, \sigma \mid \rho^3 =
\sigma^2\ \struck{\ill{}}\}$ », relié par un $=$ vertical, avec sous
l'égalité l'accolade « $\omega'^{-1} = \omega$ ». La flèche
$S\mathcal{T}_{11} \to S\Pi^{D}_{03}(4)$ porte
« $\rho \mapsto \rho\varpi^2$, $\sigma \mapsto \sigma\varpi^3$,
$\omega \mapsto \varpi^2$, $\omega^2 \mapsto \varpi^4 = \omega'$ »,
entourés d'une ligne ; la flèche à deux têtes entre $\mathcal{T}_{11}$ et
$\Pi^{D}_{03} \times \mathbb{Z}/4\mathbb{Z}$ porte
« $\rho \mapsto (\rho, 2)$, $\sigma \mapsto (\sigma, 3)$ », et à droite
« $\omega \to (1, 2)$ »}

\[
S\Pi^{D}_{03}(4) = \{\rho, \sigma, \varpi \mid
\rho^3 = \sigma^2 = \dot\varpi^{4},\ \varpi \text{ central}\}
\]
\textbf{NB} \quad $\dot\varpi^{4} = \ill{}^{\,-1} = \omega'$

\note{en haut à droite de la page. Le point au-dessus de $\varpi$ est
incertain}

\note{au bas de la page, tête-bêche :}
\[
\overbrace{S\Pi^{D}_{0\,3}(4)}^{\pi_1(U_{03}, \mathfrak{S}_3 \times \mu_4)}
\longrightarrow \Pi^{D}_{0,3}
\]

\end{document}
